\documentclass[11pt]{article}
\usepackage[a4paper,margin=25mm]{geometry}
\usepackage{amsmath,amssymb,amsthm,xcolor,hyperref}
\newtheorem{lemma}{Lemma}
\newtheorem{theorem}{Theorem}
\newcommand{\Q}{Q_n}
\title{Matching extensions: local proof audits}
\author{ARIS-assisted local derivation}
\date{November 1, 2026}
\begin{document}
\maketitle
\noindent\fcolorbox{blue!60!black}{blue!5}{\parbox{.94\linewidth}{
The result below is a restricted theorem, relative to the stated imported results.
The full six-direction problem remains open in this work. Review is same-family
and provisional; no novelty or formal-verification claim is made.}}

\section*{Statement and imports}
\begin{theorem}
Let $n\ge5$, let $M\subseteq E(\Q)$ be a matching, and let $u$ be uncovered.
Suppose a cube matching $L\supseteq M$ covers every vertex except $u,w$.
For every vertex $v$ of parity opposite to $u$, there is a Hamilton $u$--$v$
path in $\Q$ containing $M$.
\end{theorem}
The imported results are Fink's perfect-matching cycle theorem, the
prescribed-endpoint theorem of Gregor, Novotny and Skrekovski for perfect matchings
of $B(\Q)$, and the five-direction prescribed-endpoint theorem in the supplied
Fink--Hotmar paper. Their computer-assisted bases were not rerun here.

\section*{The key obstruction lemma}
\begin{lemma}
Let $P$ be a perfect matching of $B(\Q)$, let $x,y$ have opposite parity, and
suppose $xy\notin P$. Write $a=P(x)$, $b=P(y)$. If
$R=(P\setminus\{xa,yb\})\cup\{ab\}$ contains a half-layer, the cube edges of
$P$ use at most three directions.
\end{lemma}
\begin{proof}
For an $i$-half-layer with endpoint set $S$, put $T=V(\Q)\setminus S$.
Both induced cube graphs on $S,T$ have only $i$-edges. Since $R$ avoids $x,y$,
both vertices lie in $T$. If $ab$ is not in the half-layer, that entire
half-layer belongs to $P$ and saturates $S$. All short edges of $P$ then have
direction $i$. Otherwise the half-layer minus $ab$ belongs to $P$; only
$xa,yb$ can cross between $S,T$. All other short edges have direction $i$,
and the two exceptional edges introduce at most two additional directions.
\end{proof}

\section*{Main construction}
If $v=w$, add the artificial edge $uw$, apply Fink's cycle theorem, and
delete $uw$. If $uw$ is a cube edge, the known perfect-completion case
applies. If $L$ uses at most three directions, the known five-direction
endpoint theorem applies because $u$ is uncovered.

In the remaining case $uw$ is long and $L$ uses at least four directions.
Set $P=L\cup\{uw\}$. Choose a cube neighbor $x$ of $w$ other than $P(v)$;
at most one of the $n$ neighbors is excluded. Here $x\ne u$, and $x,v$
have opposite parity with $xv\notin P$.
The lemma rules out a compressed half-layer. The prescribed-endpoint theorem
gives a Hamilton $x$--$v$ path $K$ containing $P$, with all added edges short.

A Hamilton path containing a perfect matching alternates matching and
nonmatching edges: its matching edges occupy positions $(0,1),(2,3),\ldots$.
Since $x$ has the parity of $u$, the artificial edge occurs in the direction
$u\to w$. Deleting it leaves disjoint spanning components $K[x,u]$ and
$K[w,v]$. Reverse the first and add the cube edge $xw$ to obtain
\[
u\leadsto x\;--\;w\leadsto v.
\]
All vertices are visited once, all edges are now cube edges, and no edge of
$L$ was deleted. This proves the restricted theorem.

\section*{Exact remaining completion obstruction}
Let $F=V(\Q)\setminus V(M)$, and partition $F=E\cup O$, with $u\in E$ and
$E$ the parity class of $u$. The two sides have equal size. A suitable $L$
exists for some $w\in O$ exactly when
\[
|N_F(S)|\ge |S|\qquad\hbox{for every }S\subseteq E\setminus\{u\}.
\]
This is Hall's theorem on the residual graph with $u$ removed. In $Q_6$,
every $31$-edge matching is therefore covered. For a $30$-edge matching with
holes $u,z\in E$ and $w,t\in O$, the condition reduces to $z$ being adjacent
to at least one of $w,t$. Failure is a completion obstruction, not a
Hamilton-path counterexample.

\section*{Audit and evidence}
No substantive gap was found by the fresh reviewer. Two editorial items
were repaired: the dimension $n\ge2$ in the earlier cycle-deletion lemma,
and the existential quantifier on the second hole in Hall's criterion.
The dependency graph is acyclic; experiments are not proof premises.

The independent rigidity enumeration checked all $40\,320$ bipartite perfect
matchings in $Q_4$ and $2\,257\,920$ admissible endpoint pairs. All $5\,376$
half-layer cases satisfied the direction bound. Sixteen six-dimensional
instances exercised the artificial-edge deletion and splice, with every
final witness checked. Ninety-six single-bridge samples had three-bridge
witnesses; this does not prove the corresponding universal assertion.

\noindent\fcolorbox{orange!70!black}{orange!5}{\parbox{.94\linewidth}{
Still open here: matchings failing the residual Hall condition and not
covered by the other known sufficient conditions. The entire $P_6$
obligation has not been discharged.}}
\newpage
\section*{Four-hole continuation: one displaced edge}
\begin{theorem}
Let $n\ge5$, let $M$ be a cube matching avoiding $u$, and let $u,v$ have
opposite parity. Suppose a cube matching $L$ avoids exactly $u,v$ and
$|M\setminus L|\le1$. Then a Hamilton $u$--$v$ path containing $M$ exists.
\end{theorem}
\begin{proof}
If $M\subseteq L$, use the cycle-deletion argument. Otherwise write
$M\setminus L=\{ab\}$. If $L$ uses at most three directions, $M$ uses at
most four and the known five-direction endpoint theorem applies.
Otherwise set $P=L\cup\{uv\}$. This is a perfect bipartite matching with
at least four short directions. Since $u$ is uncovered by $M$, $ab\ne uv$
and $ab\notin P$. The obstruction lemma rules out a compressed half-layer
at $a,b$. The known endpoint theorem gives a Hamilton $a$--$b$ path through
$P$. Add the short edge $ab$ to close the cycle, then delete $uv$.
The result contains every edge of $M$ and has no artificial long edge.
\end{proof}
For four holes $u,z$ of one parity and $w,t$ of the other, a three-edge
augmenting path $z$--$a$--$b$--$t$ with $ab\in M$ supplies
$L=(M\setminus\{ab\})\cup\{za,bt\}$, proving the $u$--$w$ case.
This does not give every prescribed endpoint. In fact, if $v$ is covered
by $M$, the displaced edge must meet $v$ and this completion criterion
already implies the earlier root-preserving near-perfect criterion.

\section*{A necessary extra-cut bound}
For the remaining covered-target case, start with a Hamilton $x$--$v$
extension of $P=M\cup\{uw,zt\}$. Remove both artificial edges without
adding them back, remove $k$ other nonmatching edges, and add only cube
edges to obtain a $u$--$v$ path. Then
\[
k\ge\frac{\min\{d(z,w),d(z,t)\}-1}{2}.
\]
Indeed, let $D$ be the actually removed nonmatching edges and $A$ the
actually added edges. $D$ is a matching. In $A\cup D$, the excess of added
degree is one at $x,z,w,t$ and zero elsewhere. Alternating-trail
decomposition pairs these four terminals. The trail starting at $z$ must
end at $w$ or $t$, has length at most $2k+1$, and consists of cube edges.
Hamming distance gives the bound. When both distances are five, at least
two extra cuts are necessary.

\section*{Audit status and evidence}
A fresh reviewer found no substantive F1--F5 gap; the standalone dimension
assumption was made explicit. The six possible orders of three cuts were
checked independently: four reconnect to a spanning path, two leave a
separate cycle. Four new displaced-edge examples and forty covered-target
repairs passed independent witness checking. Sixteen of the latter have
distance-five holes and attain the two-extra-cut lower bound.

\noindent\fcolorbox{orange!70!black}{orange!5}{\parbox{.94\linewidth}{
Still unproved: universal existence of compatible extra cuts for every
remaining matching and covered target. The forty examples are not an
exhaustive proof. Review remains same-family and provisional.}}
\end{document}
